
\documentclass[11pt]{article}
\usepackage[margin=1in]{geometry}
\usepackage{amsmath,amssymb,mathtools}
\usepackage{booktabs,longtable,array}
\usepackage{hyperref}
\usepackage{enumitem}
\usepackage{natbib}
\usepackage{caption}
\usepackage{microtype}
\numberwithin{equation}{section}
\newtagform{brackets}{[}{]}
\usetagform{brackets}
\hypersetup{hidelinks}
\newcommand{\IRFR}{\mathrm{IRFR}}
\newcommand{\MSNR}{\mathrm{MSNR}}
\newcommand{\dd}{\,\mathrm{d}}
\title{\textbf{Empirical Evidence for a Unified Boundary-Field Theory of Interest Rates}\\\large Yield-Curve Shape, Volatility, and Risk-Neutral Structure from a Single Dataset}
\author{Working paper draft}
\date{Updated equation-numbered PDF review version: 20 June 2026}
\begin{document}
\maketitle
\begin{abstract}
This paper presents empirical evidence for a Unified Boundary-Field Theory of Interest Rates, developed using the instantaneous rolling forward rate (IRFR) as its primitive. The central empirical question is whether a single maturity-space boundary geometry can organise multiple observables extracted from the same interest-rate dataset: yield-curve shape, the spatial decay parameter $a^2$, interest-rate volatility, covariance rank, and market-price-of-risk structure.

Using U.S. fitted instantaneous forward rates at benchmark maturities 1Y, 3Y, 5Y, 10Y, and 30Y, the paper estimates the boundary-field state implied by the IRFR framework. The short end is represented by a policy-sensitive short boundary and curvature state, while the long end is represented by a stochastic long-maturity boundary. The maturity field between them is governed by the spatial decay parameter $a^2$. The empirical tests ask whether the same estimated geometry that fits curve shape also generates the observed maturity-volatility structure and supports a risk-neutral drift restriction.

The results are supportive. The benchmark specification identifies a stable long-sample spatial decay close to $a^2=0.09$. The covariance evidence favours a two-boundary-risk structure, separating short/curvature risk from long-boundary risk. A one-Brownian specification can reproduce some volatility ordering, but it fails to match volatility scale and performs materially worse in market-signal-to-noise-ratio diagnostics of the market-price-of-risk restriction. Fixed $a^2$ generates strong maturity-volatility geometry over structural windows, especially over rolling twenty-year samples. The market-price-of-risk restriction performs strongly against a zero-drift benchmark, while centred realised-drift tests are unstable because cross-sectional realised drift is itself a fragile empirical object.

The 30-year fitted forward is used as the benchmark observable proxy for the stochastic long boundary. The 20-year point is examined as an alternative effective long-boundary proxy. Although the 20-year specification gives a similar full-sample estimate of $a^2$, it produces less stable period and rolling-window volatility correlations. The 30-year point is therefore retained as the benchmark long-boundary proxy, while the 20-year evidence is interpreted as confirming the difficulty of measuring an asymptotic boundary from finite maturities.
\end{abstract}

\section{Introduction}
Interest-rate data are usually analysed through separate empirical lenses. Yield-curve shape is studied through curve-fitting models or factor decompositions. Volatility is studied through covariance matrices, principal components, or term-structure volatility models. Risk-neutral structure is studied through pricing kernels, affine restrictions, or market prices of risk. These objects are closely related in theory, but they are often estimated and evaluated separately.

This paper asks whether they can be unified empirically. More specifically, it asks whether the empirical data support the IRFR claim that yield-curve shape and volatility structure can be organised within a single maturity-space boundary framework. The test is not whether the boundary-field theory can be proved from historical data alone. Rather, the question is whether the main observables extracted from one interest-rate dataset are jointly consistent with the theory's central implication: that the same boundary geometry should organise curve shape, maturity volatility, covariance structure, and the admissible market-price-of-risk restriction.

The hypothesis is that interest rates form a maturity-space boundary field. The short end of the curve is a policy-sensitive boundary. The long end is a stochastic boundary reflecting the market's long-run nominal-rate anchor. Between them lies a maturity field governed by a spatial decay parameter $a^2$. If this representation is more than a curve-fitting device, the same boundary geometry should appear in several observables extracted from the same dataset.

\section{Boundary-field state representation}
At date $t$, the forward-rate field at maturity $\tau$ is represented by
\begin{equation}
 x_t(\tau)=X_t+\left[(r_t-X_t)+B_t\tau\right]e^{-a^2\tau}.
 \label{eq:state}
\end{equation}
Here $r_t$ is the short boundary, $B_t$ is the curvature or policy-deformation state, $X_t$ is the stochastic long boundary, and $a^2$ controls the spatial decay of short-boundary effects into maturity space. The associated maturity-loading vector is
\begin{equation}
 J(\tau)=\left(e^{-a^2\tau},\;\tau e^{-a^2\tau},\;1-e^{-a^2\tau}\right).
 \label{eq:J}
\end{equation}
This vector maps state increments into forward-rate increments:
\begin{equation}
 \dd x_t(\tau)=e^{-a^2\tau}\dd r_t+\tau e^{-a^2\tau}\dd B_t+(1-e^{-a^2\tau})\dd X_t.
 \label{eq:dx}
\end{equation}
The same loading vector is used throughout the empirical exercise. It is not re-estimated separately for shape, volatility, covariance, and risk-price tests.

\section{Dataset and empirical construction}
The benchmark dataset uses U.S. fitted instantaneous forward rates at maturities
\begin{equation}
 \mathcal{T}=\{1,3,5,10,30\}.
 \label{eq:benchmark-tenors}
\end{equation}
The fitted-forward columns are the Federal Reserve Board nominal yield-curve series $\texttt{SVENF01}$ through $\texttt{SVENF30}$, based on the Gürkaynak--Sack--Wright U.S. Treasury yield-curve dataset \citep{GurkaynakSackWright2007,FederalReserveYieldCurve}. For each observation date, the empirical forward-rate vector is
\begin{equation}
 F_t=\big(f_t(1),f_t(3),f_t(5),f_t(10),f_t(30)\big).
 \label{eq:empirical-vector}
\end{equation}
The 30-year point is used as the observable proxy for the stochastic long boundary,
\begin{equation}
 B_t^L := f_t(30).
 \label{eq:long-boundary-proxy}
\end{equation}
The 20-year point is excluded from the benchmark and retained as a robustness exercise because it behaves less consistently as a long-boundary proxy.

For a given value of $a^2$, the observed benchmark forward curve is fitted to equation~\eqref{eq:state}. This produces a time series of estimated states $(r_t,B_t,X_t)$. The spatial decay parameter $a^2$ is estimated by cross-sectional shape fit, using the least-squares criterion
\begin{equation}
 \widehat{a^2}=\arg\min_{a^2}\sum_t\sum_{\tau\in\mathcal{T}}\left[f_t(\tau)-x_t(\tau;a^2)\right]^2.
 \label{eq:a2-objective}
\end{equation}
In the benchmark specification, the full-sample estimate is close to $a^2=0.091893$.

\section{Boundary-risk decomposition}
The empirical covariance structure is constructed from increments of the estimated boundary states. The theory separates the short and curvature states from the long-boundary state:
\begin{equation}
 \dd r_t=\nu_r\dd W_t^S,\qquad \dd B_t=\nu_B\dd W_t^S,\qquad \dd X_t=\nu_X\dd W_t^X.
 \label{eq:risk-states}
\end{equation}
The two Brownian directions are allowed to be correlated according to
\begin{equation}
 \dd W_t^S\dd W_t^X=\rho\dd t.
 \label{eq:brownian-correlation}
\end{equation}
The coefficients $\nu_r$, $\nu_B$, and $\nu_X$ are signed state-level volatility loadings. The short/curvature Brownian direction $W_t^S$ is identified from the first principal component of $(\dd r_t,\dd B_t)$, while the long-boundary direction is estimated from $\dd X_t$.

The corresponding maturity-risk loadings are
\begin{align}
 g_S(\tau)&=e^{-a^2\tau}\nu_r+\tau e^{-a^2\tau}\nu_B,
 \label{eq:short-boundary-loading}\\
 g_X(\tau)&=(1-e^{-a^2\tau})\nu_X.
 \label{eq:long-boundary-loading}
\end{align}
The maturity variance is
\begin{equation}
 g_{\mathrm{eff}}^2(\tau)=g_S^2(\tau)+g_X^2(\tau)+2\rho g_S(\tau)g_X(\tau).
 \label{eq:effective-variance}
\end{equation}
Equivalently, if $L=\begin{pmatrix}\nu_r&0\\\nu_B&0\\0&\nu_X\end{pmatrix}$ and $C=\begin{pmatrix}1&\rho\\\rho&1\end{pmatrix}$, then the model-implied maturity covariance matrix is
\begin{equation}
 K=JLC L^\top J^\top.
 \label{eq:covariance-kernel}
\end{equation}
This covariance matrix has rank at most two. It becomes rank one only if the two Brownian directions collapse into a single effective Brownian motion.

\section{Market-price-of-risk and MSNR diagnostics}
The same two maturity-risk directions imply a market-price-of-risk restriction:
\begin{equation}
 f(\tau)=g_S(\tau)\lambda_S+g_X(\tau)\lambda_X,
 \label{eq:mpr}
\end{equation}
where $f(\tau)$ is the estimated physical drift of the forward-rate increment and $\lambda_S,\lambda_X$ are the market prices of short/curvature risk and long-boundary risk.

MSNR is not used as a synonym for the market price of risk. The market prices of risk are $\lambda_S$ and $\lambda_X$ in equation~\eqref{eq:mpr}. MSNR is the empirical diagnostic used to test how well the realised drift vector is explained by the boundary-risk projection. Two benchmarks are reported:
\begin{align}
 R^2_{\mathrm{zero}}&=1-\frac{\sum_\tau(f_\tau-\widehat f_\tau)^2}{\sum_\tau f_\tau^2},
 \label{eq:zero-msnr}\\
 R^2_{\mathrm{centred}}&=1-\frac{\sum_\tau(f_\tau-\widehat f_\tau)^2}{\sum_\tau(f_\tau-\bar f)^2}.
 \label{eq:centred-msnr}
\end{align}
The zero benchmark asks whether the model explains realised drift relative to no drift. The centred benchmark asks whether the model explains the demeaned cross-maturity drift shape. The latter is stricter and less stable because realised drift is noisy and the centred denominator can be very small.

\section{Empirical results}
\subsection{Branch selection}
Table~\ref{tab:branch} reports the branch diagnostic. Shape fit is high for all branches, so shape alone does not identify $r=1$. The volatility-transition and scale diagnostics, however, favour the repeated-root branch. As $r$ moves toward one, transition-region volatility RMSE falls monotonically from 252.5bp to 68.4bp, and the median model/actual volatility ratio moves from 1.55 toward 1.10.

\begin{table}[htbp]\centering\caption{Branch diagnostic}\label{tab:branch}\begin{tabular}{rrrrrr}\toprule
$r$ & Shape $R^2$ & Vol. corr. & 5Y/10Y/30Y RMSE bp & Median ratio & Zero MSNR $R^2$\\\midrule
0.50&0.9934&0.9304&252.5&1.55&0.8545\\
0.60&0.9938&0.9263&185.0&1.39&0.8569\\
0.70&0.9942&0.9222&139.3&1.28&0.8552\\
0.80&0.9945&0.9181&107.5&1.20&0.8514\\
0.90&0.9946&0.9140&84.9&1.14&0.8467\\
1.00&0.9947&0.9102&68.4&1.10&0.8421\\\bottomrule\end{tabular}\end{table}

\subsection{One Brownian versus two Brownian directions}
Table~\ref{tab:brownianfull} compares one-Brownian and two-Brownian specifications over the full sample. The one-Brownian model has a higher raw volatility correlation, but this is misleading because it does materially worse on volatility scale and risk-price projection. The two-Brownian model reduces full-sample volatility RMSE from 91.7bp to 53.5bp and improves zero-benchmark MSNR from 0.261 to 0.842.

\begin{table}[htbp]\centering\caption{One-Brownian vs two-Brownian full-sample comparison}\label{tab:brownianfull}\begin{tabular}{lrr}\toprule
Diagnostic&One Brownian&Two Brownian\\\midrule
Volatility correlation&0.988&0.910\\
Volatility RMSE, bp&91.7&53.5\\
Zero-benchmark MSNR $R^2$&0.261&0.842\\
Centred MSNR $R^2$&-76.76&-15.61\\\bottomrule\end{tabular}\end{table}

Table~\ref{tab:brownianperiods} shows the period-level comparison. The most robust evidence in favour of two Brownian factors is the MSNR/risk-price test: the two-Brownian model improves zero-benchmark MSNR in every period.

\begin{table}[htbp]\centering\caption{One-Brownian vs two-Brownian period comparison}\label{tab:brownianperiods}\begin{tabular}{lrrrrrr}\toprule
Period&One corr.&Two corr.&One RMSE&Two RMSE&One zero MSNR&Two zero MSNR\\\midrule
1980--1990&0.989&0.957&90.6&195.6&0.263&0.892\\
1990--2000&0.990&0.932&75.5&56.3&0.471&0.846\\
2000--2010&0.966&0.833&82.8&53.2&0.049&0.956\\
2010--2020&0.549&0.072&59.3&43.3&0.649&0.721\\
2020--2025&0.559&0.846&91.1&27.4&0.647&0.976\\
Full sample&0.988&0.910&91.7&53.5&0.261&0.842\\\bottomrule\end{tabular}\end{table}

\subsection{Spatial decay stability}
Table~\ref{tab:a2periods} reports the spatial-decay estimates. The clean long-sample estimate is $a^2\approx0.0919$. Shorter windows reveal regime sensitivity, especially after 2020, but the structural-horizon evidence supports a stable decay scale around 0.09.

\begin{table}[htbp]\centering\caption{Spatial-decay estimates}\label{tab:a2periods}\begin{tabular}{lrrrrr}\toprule
Window&Obs.&$a^2_{\mathrm{shape}}$&$1/(2a^2)$&Shape $R^2$&RMSE bp\\\midrule
1980--1990&1025&0.0493&10.13&0.9851&12.04\\
1990--2000&2483&0.0967&5.17&0.9800&18.27\\
2000--2010&2086&0.0825&6.06&0.9802&17.21\\
2010--2020&1146&0.1300&3.85&0.9840&17.11\\
2020--2025&867&0.2352&2.13&0.9953&9.67\\
Full sample&7607&0.0919&5.44&0.9947&17.06\\\bottomrule\end{tabular}\end{table}

\subsection{Fixed-$a^2$ volatility geometry}
With full-sample $a^2$ fixed at 0.091893, the model continues to produce strong maturity-volatility geometry over the full sample and rolling twenty-year windows. Table~\ref{tab:test4} reports the period results. The 2010--2020 period is the main weak decade, with volatility correlation 0.0332, consistent with regime-specific ZIRP/QE distortion in the realised volatility shape.

\begin{table}[htbp]\centering\caption{Volatility from fixed full-sample $a^2$}\label{tab:test4}\begin{tabular}{lrrrrrrr}\toprule
Window&Obs.&Shape $R^2$&Vol corr.&RMSE bp&Median ratio&Within 20\%&$\rho_{SX}$\\\midrule
1980--1990&1025&0.9805&0.9495&49.4&0.989&2&0.6869\\
1990--2000&2483&0.9800&0.9304&62.5&1.201&2&0.7451\\
2000--2010&2086&0.9799&0.8325&43.0&1.090&3&0.3810\\
2010--2020&1146&0.9816&0.0332&42.1&0.711&2&0.1900\\
2020--2025&867&0.9899&0.8596&26.1&0.857&4&0.3097\\
Full sample&7607&0.9947&0.9102&53.5&1.096&3&0.5365\\\bottomrule\end{tabular}\end{table}

The corrected rolling-window evidence is stronger than the decade comparison alone suggests: rolling twenty-year fixed-$a^2$ windows have median volatility correlation 0.8908 and minimum volatility correlation 0.8681 across 21 windows. This supports the interpretation that the geometry is stable over structural horizons even when shorter realised windows remain regime-sensitive.

\subsection{Market-price-of-risk diagnostics}
Table~\ref{tab:msnrperiods} reports MSNR results. The zero-drift benchmark is consistently strong, while the centred benchmark is unstable and can become severely negative when the centred denominator is tiny. The full-sample denominator ratio is only 0.0095, and the 2020--2025 denominator ratio is only 0.0058. The centred statistic is therefore asking the model to explain a very small demeaned cross-maturity drift shape.

\begin{table}[htbp]\centering\caption{MSNR period results, fixed $a^2$}\label{tab:msnrperiods}\begin{tabular}{lrrrrrr}\toprule
Window&Centred $R^2$&Zero $R^2$&Denom. ratio&$\lambda_S$&$\lambda_X$&$\rho_{SX}$\\\midrule
1980--1990&0.6536&0.9478&0.1508&0.0146&-0.2201&0.6869\\
1990--2000&-0.3980&0.8477&0.1090&-0.0528&-0.0692&0.7451\\
2000--2010&0.8185&0.9465&0.2948&-0.5505&-0.1660&0.3810\\
2010--2020&0.1127&0.6559&0.3877&-0.0946&-0.5137&0.1900\\
2020--2025&-7.4128&0.9509&0.0058&0.3871&0.4430&0.3097\\
Full sample&-15.6078&0.8421&0.0095&-0.1072&-0.1070&0.5365\\\bottomrule\end{tabular}\end{table}

The conclusion is not that centred MSNR disproves the boundary geometry. Rather, the model prices drift well relative to zero drift, but not reliably relative to a sample-mean cross-sectional drift benchmark. This is exactly what one should expect when realised mean drift is a fragile empirical object.

\section{Long-boundary proxy and robustness}
A natural limitation of the empirical test is that the stochastic long boundary is not directly observable. In the benchmark specification, the 30-year fitted forward is used as the observable proxy for long-boundary variation. This choice is empirically useful but imperfect. The 30-year fitted-forward node may contain extrapolation, liquidity, issuance, pension/LDI, or other long-end market-structure effects.

The 20-year point is tested as an alternative effective long-boundary proxy. Although it gives a similar full-sample estimate of $a^2$, it produces less stable period and rolling-window volatility geometry. This suggests that 20-year data sometimes behave like a transition-region maturity rather than a clean long-boundary proxy. The 30-year point is therefore retained as the benchmark, while the 20-year evidence is treated as a robustness warning about finite-maturity measurement of an asymptotic boundary.

\section{Conclusion}
Taken together, the evidence suggests that the same boundary-field geometry which organises yield-curve shape also organises volatility and risk-neutral drift restrictions. The evidence does not imply that interest-rate risk premia are globally stationary, that realised mean drift is easy to estimate, or that every maturity is measured without error. It supports a more specific claim: interest-rate observables contain a common boundary-field structure linking shape, volatility, covariance, and risk-neutral character.

The strongest empirical point is cumulative. The model is not judged by one statistic. It is judged by whether one estimated boundary geometry supports branch selection, spatial decay stability, two-boundary covariance, fixed-$a^2$ maturity-volatility geometry, and market-price-of-risk projection. The results are supportive under this linked diagnostic architecture, but they remain empirical evidence rather than proof of uniqueness or canonicality.

\appendix
\section{Replication package}
The accompanying replication package contains:
\begin{itemize}[leftmargin=2em]
\item \texttt{data/derived/}: CSV files for all reported diagnostic tables;
\item \texttt{data/raw/README\_raw\_data.md}: raw-data source and required Fed columns;
\item \texttt{code/replicate\_ubft\_empirical.py}: transparent Python replication script;
\item \texttt{code/requirements.txt}: Python dependencies;
\item \texttt{paper/references.bib}: bibliography used by this LaTeX source.
\end{itemize}


\begin{thebibliography}{2}

\bibitem[G{\"u}rkaynak et~al.(2007)]{GurkaynakSackWright2007}
G{\"u}rkaynak, R. S., Sack, B., and Wright, J. H. (2007).
\newblock The U.S. Treasury Yield Curve: 1961 to the Present.
\newblock \emph{Journal of Monetary Economics}, 54(8), 2291--2304.

\bibitem[Federal Reserve Board(2026)]{FederalReserveYieldCurve}
Board of Governors of the Federal Reserve System (2026).
\newblock Nominal Yield Curve: G{\"u}rkaynak--Sack--Wright Treasury Yield Curve Estimates.
\newblock Federal Reserve Board yield-curve data tables. Data columns include \texttt{SVENF01}--\texttt{SVENF30} fitted instantaneous forward rates.

\end{thebibliography}

\end{document}
